% SUBMISSION: {"kind":"research_attempt","category":"erdos","problem_id":"757","model":{"provider":"OpenAI","version":"Codex","version_uncertainty":"The retained proof-run record exposes Codex but does not record its exact underlying checkpoint; current submission preparation identifies GPT-6-based Codex without establishing the earlier run checkpoint.","reasoning":"Not exposed","generated_on":"2026-10-15"},"other_models":"None documented; no additional model identity inferred.","tools":"Lean 4.24.0, pinned mathlib, Python exact arithmetic, GitHub Actions, Tectonic 0.17.0.","target":"Determine the optimal c such that every finite real (4,5)-set B has a Sidon subset of size at least c times |B|.","scope":"Every weak Sidon 16-point real set containing the specified Ma-Tang 14-point block has a Sidon 10-subset. No new global bound or full solution.","human_contributions":"Requester set research objectives and publication constraints; salemon maintains and submits artifacts. No independent human mathematical contribution or checking recorded.","verification":"Lean build and transitive axiom audit passed all 52 registered project declarations; exact certificate checker passed. Current scope audit is model-assisted, not independent human review.","reproducibility":"Complete pinned paper, code and data at salemon/erdos-242 commit ca3f08715aeffb98581410a67bff9d02fb9ecbfe. Prompt A/B criteria applied retrospectively to this description, not claimed as original proof-run prompts. Private transcript not supplied."}
% ATTEMPT_STATUS: partial
% FIRST_POSTED: 2026-10-15
\documentclass[11pt]{article}
\usepackage[T1]{fontenc}
\usepackage{lmodern,amsmath,amssymb}
\usepackage[margin=1in]{geometry}
\usepackage{xurl}
\usepackage[colorlinks=true,urlcolor=blue]{hyperref}
\title{Erd\H{o}s \#757: a Lean-checked two-point extension obstruction}
\author{OpenAI Codex (proof-run model version unrecorded)}
\date{15 October 2026}
\begin{document}
\maketitle

\section{Original problem and exact restricted statement}
Erd\H{o}s problem \#757 asks for the optimal constant $c_*>0$ such that
every finite $(4,5)$-set $B\subset\mathbb R$ contains a Sidon subset
of size at least $c_*|B|$~\cite{erdos,source}. Here a $(4,5)$-set means
that each four-element subset determines at least five different positive
distances. A set is Sidon if the sums $a+b$, indexed by $a\leq b$, are
distinct, including doubled elements. Weak Sidonicity restricts that
comparison to $a<b$.

This submission describes a formalization of the following restricted theorem.
Let
\[
\begin{split}
A_0=\{&0,136,200,243,246,249,272,286,\\
&298,323,400,528,596,1056\}.
\end{split}
\]
For every finite weak Sidon set $B\subset\mathbb R$ satisfying
$A_0\subseteq B$ and $|B|=16$, there exists a Sidon subset $T\subseteq B$
with $|T|=10$. The result also holds with the $(4,5)$ hypothesis in place
of weak Sidonicity. The two added points may be arbitrary real numbers.

This rules out obtaining a sixteen-point example with largest Sidon subset
of size at most nine by appending two points to this exact block.
It does not determine $c_*$, improve its published bounds, or address all
sixteen-point sets. Replacing old points or adding more than two points
is outside the theorem's scope.

\section{Proof mechanism, literature, and complete output}
Ma--Tang's block and their bounds $9/17\leq c_*\leq4/7$ are existing
results~\cite{matang}. The standard facts that a weak Sidon $n$-set
has at most $n-2$ nontrivial three-term arithmetic progressions for $n\geq2$,
and that its Sidon subsets are precisely its progression-free subsets,
are reproved in the accompanying Lean development. The implication
from $(4,5)$ to weak Sidonicity is also proved there.

The exact block has twelve progressions. An explicit twelve-row certificate
shows that deleting any at most two of its points leaves a Sidon eight-set.
A sixteen-point weak Sidon extension therefore has at most two new
progressions. Each contains an old point, since only two points were added.
Choose one old point from each new progression. By the certificate, select
a Sidon eight-set avoiding the chosen points, and adjoin the two new points.
No old progression survives in that union, and every new progression was
hit by a deleted old point. The union is consequently Sidon and has size ten.

The complete output, definitions, finite table, proof of the general extension
lemma, and limitations are available in the
\href{https://github.com/salemon/erdos-242/blob/ca3f08715aeffb98581410a67bff9d02fb9ecbfe/paper/erdos757.pdf}{six-page paper (PDF)}
and its
\href{https://github.com/salemon/erdos-242/blob/ca3f08715aeffb98581410a67bff9d02fb9ecbfe/paper/erdos757.tex}{LaTeX source}.
The repository name reflects an earlier \#242 baseline; this theorem is
about \#757.

The original investigation inspected Ma--Tang v1; paper preparation additionally
checked v2, including Proposition 2.2, Lemmas 2.3--2.4 and 5.1.
Prior public notes already report one-point and deletion observations, as
credited in the
\href{https://github.com/salemon/erdos-242/blob/ca3f08715aeffb98581410a67bff9d02fb9ecbfe/docs/757/REPORT.md}{research report}.
Neither the block nor those observations are claimed as discoveries here.
No exact antecedent for the two-point argument was found in that limited
review, but mathematical novelty, significance, and formalization priority
remain unconfirmed. The observation may be a straightforward known corollary.
Existing \#757 manuscripts in this collection are preserved.

\section{Lean source and verification scope}
The
\href{https://github.com/salemon/erdos-242/tree/ca3f08715aeffb98581410a67bff9d02fb9ecbfe/Erdos757}{complete Lean source}
contains the definitions, general proof, finite certificates, and
integer-to-real transfer. In \texttt{Block.lean}, the relevant declarations are
\path{Erdos757.ma_tang_two_point_extension} and
\path{Erdos757.ma_tang_four_five_extension}.
Their concrete AP-count and resilience hypotheses are discharged in Lean;
they are not left as assumptions supplied by a Python program.

Lean 4.24.0 and mathlib commit
\path{f897ebcf72cd16f89ab4577d0c826cd14afaafc7} are pinned.
On 15 October 2026, \texttt{lake build} and the project's transitive axiom
audit passed for all 52 registered declarations across both developments.
Only the standard \texttt{propext}, \texttt{Classical.choice}, and
\texttt{Quot.sound} are permitted. There are no \texttt{sorry}/\texttt{admit}
placeholders, additional mathematical axioms, or \texttt{native\_decide}
in the proof source. Kernel-reduced \texttt{decide} checks the finite witnesses.
The separate exact Python checker passed 91 deletion cases, counted twelve
progressions and 143 Sidon eight-subsets, and checked the block's $(4,5)$
property. The general theorem is not merely those finite checks.

Reproduction instructions and scripts are in the
\href{https://github.com/salemon/erdos-242/blob/ca3f08715aeffb98581410a67bff9d02fb9ecbfe/README.md}{pinned README}.
The formal result does not include a proof of the optimal global constant.
No independent human proof review is recorded, and no verification by this
collection's maintainers is implied. This is a multi-file Lean project;
no tested single-file Lean Web reproduction is supplied.

\section{Model, contributions, and process}
Provider/tool: OpenAI Codex. The retained proof-run artifacts identify Codex
but do not record an exact underlying checkpoint or reasoning configuration;
those details cannot be independently recovered from this submission and are
not invented. The current submission assistant is identified as GPT-6-based
Codex; this does not establish the exact checkpoint of the earlier proof run.
No other material model is documented. Generation and submission preparation
date: 15 October 2026.

Codex performed literature searches, target selection, proof development,
Lean/Python implementation, automated checks, and exposition. The human
requester set objectives and constraints and requested publication;
\texttt{salemon} is the repository maintainer/submitter, not a claimed
independent mathematical reviewer. No independently checked human mathematical
contribution or endorsement is asserted. Tools used include Lean/mathlib,
Python exact arithmetic, GitHub Actions, and Tectonic for PDF preparation.

The original research predates this collection's prompt workflow. For this
description, the preparation and adversarial-audit criteria in prompts A and B
were applied retrospectively: retain the full original target, trace the proof
dependencies, distinguish the restricted theorem from the global problem,
check the actual axiom report, and preserve the unresolved status. Those prompts
were not used as the initial prompts of the historical proof run.
The original private chat transcript is not supplied. Public code, data,
source-check records, paper, and verification reports are linked above.
The present audit is model-assisted, not independent human checking.

\section{Completion Estimate}
\noindent\textbf{COMPLETION: 0\%}
This is a deliberately conservative estimate toward determining the original
optimal constant: no improvement in its global interval is claimed.
The proved local obstruction excludes one construction strategy but provides
no defensible numerical measure of how much of the full problem is solved.
The value is not a confidence estimate for the restricted Lean theorem.

\section{Final status}
PARTIAL: the stated two-point obstruction is formally proved. The original
problem remains unresolved in this work. Novelty and significance are
unconfirmed; no new global bound or complete solution is claimed.

\begin{thebibliography}{9}
\bibitem{erdos} P. Erd\H{o}s, \emph{Some old and new problems in various
branches of combinatorics}, Discrete Mathematics 165/166 (1997), 227--231,
p.~231; bibliographic identification as cited by Ma--Tang.
\bibitem{source} T. F. Bloom, \emph{Erd\H{o}s Problem \#757},
\url{https://www.erdosproblems.com/757}.
\bibitem{matang} J. Ma and Q. Tang, \emph{Largest Sidon subsets in weak
Sidon sets}, arXiv:2602.23282v2, 6 March 2026,
\url{https://arxiv.org/abs/2602.23282v2}.
\end{thebibliography}
\end{document}
